\section{From formulas to differential systems}
\label{sec:formula-systems}

The word functions from Section~\ref{sec:iterated-integrals} turn a nonzero
polynomial $f=\sum_w f_wx_w$ into the nonzero series
$h_f=\sum_w f_wh_w$.  We now express this series as a scalar output of a
differential system whose dimension is at most twice the formula size.
This small system will control how late the first nonzero coefficient of
$h_f$ can occur. The underlying formula-to-path construction is due to
Nisan~\cite[Section~2, Lemma~1]{Nisan1991}; see also Raz and
Shpilka~\cite[Section~2.1, Lemma~2, Step~1]{RazShpilka2005}.

\begin{lemma}[Acyclic coefficient representation]
\label{lem:formula-representation}
Let $F$ be a field, and let $f\in F\langle x_1,\ldots,x_n\rangle$
be computed by a fan-in-two formula with at most $s$ gates, counting
leaves.  There are an integer $m\le 2s$, a row
$u\in F^{1\times m}$, a column $v\in F^{m\times1}$, and strictly
upper triangular matrices $B_1,\ldots,B_n\in F^{m\times m}$ such that
\begin{equation}
\label{eq:formula-representation}
 f_w=uB_wv\quad\text{for every word }w,
 \qquad
 B_{i_1\cdots i_k}=B_{i_1}\cdots B_{i_k},\quad B_\epsilon=I_m.
\end{equation}
In particular, $B_w=0$ whenever $|w|\ge m$.
\end{lemma}

\begin{proof}
Construct a directed acyclic graph for each subformula, with a designated
source having no incoming edges and a designated sink having no outgoing
edges.  Each edge is labeled either by a variable or by a scalar in $F$.
The label of a path is the ordered product of its edge labels.

A leaf gives two vertices joined by one edge with the leaf's label.
For a sum, take disjoint copies of the two child graphs, add a new source
and sink, and join the new source to both child sources and both child
sinks to the new sink by edges labeled $1$.  For an ordered product,
take disjoint child graphs and join the first child's sink to the second
child's source by an edge labeled $1$; retain the first source and the
second sink.  These operations preserve acyclicity and the stated
source and sink properties.  A path through a sum chooses one child,
whereas a path through a product traverses the first child and then the
second.  Induction therefore shows that the sum of source-to-sink path
labels is $f$.

If the formula has $L$ leaves and $a$ addition gates, the graph has
$m=2L+2a\le2s$ vertices: multiplication adds no vertex.  Order the
vertices topologically.  Let $E$ be the adjacency matrix of the scalar
edges, with their scalar weights, and let $N_i$ be the adjacency matrix
of the edges labeled $x_i$, with weight $1$.  In each adjacency matrix,
row $a$ and column $b$ record edges from $a$ to $b$; parallel-edge weights
are summed.  All these matrices are strictly upper triangular.  The matrix
\[
 Q=(I_m-E)^{-1}=I_m+E+\cdots+E^{m-1}
\]
records all paths consisting only of scalar edges, including the empty
path at each vertex.  Write $e_j$ for the $j$th standard column vector.
If $\sigma$ and $\rho$ are the source and sink,
every path with variable word $i_1\cdots i_k$ decomposes uniquely into
scalar subpaths separated by those $k$ variable edges.  Consequently
\[
 f_{i_1\cdots i_k}
 =e_\sigma^{\mathsf T}Q N_{i_1}Q\cdots N_{i_k}Qe_\rho,
 \qquad
 f_\epsilon=e_\sigma^{\mathsf T}Qe_\rho.
\]
Set
\[
 u=e_\sigma^{\mathsf T}Q,\qquad v=e_\rho,
 \qquad B_i=N_iQ.
\]
Then \eqref{eq:formula-representation} follows.  Multiplying a strictly
upper triangular matrix by an upper triangular matrix preserves strict
upper triangularity, so each $B_i$ has that property.  A product of $m$
strictly upper triangular $m\times m$ matrices is zero, proving the
last assertion.
\end{proof}

This construction allows arbitrary scalar leaves and constant terms.
Its size bound uses the formula tree: the child graphs represent
disjoint subtrees, so each gate is counted once.  The construction does
not enumerate the monomials of $f$.

\begin{proposition}[The system associated to a formula]
\label{prop:formula-system}
Let $F$ have characteristic zero, and let
$0\ne f\in F\langle x_1,\ldots,x_n\rangle$ have a formula of size at
most $s$.  There are $m\le2s$, a nonzero constant row $u$, a constant
column $v$, and an invertible matrix $G\in F[[z]]^{m\times m}$ such that
\begin{equation}
\label{eq:formula-system}
 G'=AG,\qquad G(0)=I_m,\qquad
 A(z)=\frac{P(z)}{q(z)},\qquad
 q(z)=\prod_{i=1}^n(z+i),\quad \deg P\le n-1,
\end{equation}
where $P$ is a polynomial matrix.  Moreover,
\begin{equation}
\label{eq:formula-output}
 h_f(z)=\sum_w f_wh_w(z)=uG(z)v\ne0.
\end{equation}
\end{proposition}

\begin{proof}
Take the representation from Lemma~\ref{lem:formula-representation}
and define the finite sum
\[
 G(z)=\sum_{|w|<m}B_wh_w(z),\qquad
 A(z)=\sum_{i=1}^n\frac{B_i}{z+i}.
\]
The empty word contributes $I_m$, and every other $h_w$ has zero
constant term.  Thus $G(0)=I_m$, which makes $G$ invertible over
$F[[z]]$.  The identity $h_{iw}'=h_w/(z+i)$ gives
\[
 G'
 =\sum_{i=1}^n\sum_{|w|<m-1}
       \frac{B_iB_wh_w}{z+i}
 =AG.
\]
The terms with $|w|=m-1$ added in the last equality vanish because
$B_iB_w=0$.  In particular the system acts on the \emph{left}, in
the same order as the letters of a word.

The denominator in \eqref{eq:formula-system} is nonzero at zero:
$q(0)=n!\ne0$.  Its numerator is
\[
 P(z)=\sum_{i=1}^n B_i\prod_{j\ne i}(z+j),
\]
whose entries have degree at most $n-1$.  Finally,
\eqref{eq:formula-representation} gives $uGv=h_f$.
This series is nonzero by Lemma~\ref{lem:word-independence}, since
$f\ne0$; in particular $u\ne0$.
\end{proof}

We have obtained a nonzero series $h_f$ inside a system of dimension at
most $2s$, with a denominator of degree $n$ and a regular expansion
point at zero.  The next section bounds its vanishing order using only
these parameters, uniformly in all scalar coefficients of the formula.
